\documentclass[12pt]{report}

\usepackage{setspace}
\renewcommand{\baselinestretch}{1.5}
\usepackage{float}
\usepackage{lipsum}
\usepackage{graphicx}
\usepackage{subcaption}
\usepackage[table,xcdraw,dvipsnames]{xcolor}
\usepackage{titlesec}
\usepackage{titletoc}
\usepackage{titling}
\usepackage[T1]{fontenc}
\usepackage[utf8x]{inputenc}
\usepackage{amsmath}
\usepackage{caption}
%\usepackage[a4paper,margin=0.5in]{geometry}
\usepackage[a4paper,%
            left=2cm,right=2cm,top=1cm,bottom=1cm,%
            footskip=.25in]{geometry}
\usepackage[normalem]{ulem}
\usepackage[none]{hyphenat}
\usepackage[hidelinks]{hyperref}
\PassOptionsToPackage{hyphens}{url}\usepackage{hyperref}
\useunder{\uline}{\ul}{}

% No hyphenation (word wrapping) with text justification
\tolerance=1
\emergencystretch=\maxdimen
\hyphenpenalty=10000
\hbadness=10000

% Replace table caption name
\usepackage{caption}
\captionsetup[table]{name=Tableau}

% For toc, lot and lof
\usepackage[titles]{tocloft}

% For liste of figures and tables!
\graphicspath{{figures/}}
\usepackage{array}

% Disable auto indentation globally
\setlength{\parindent}{0pt}

% Font options
\usepackage{mathptmx}
%\usepackage{tgtermes}
%\usepackage{tgpagella}
%\usepackage{helvet}
%\usepackage{tgbonum}

% fancy logo for page numbers!
\usepackage{fancyhdr,blindtext,tikz}
\usepackage{lastpage,refcount}

% Box frame for figures #\fbox{}
\setlength{\fboxsep}{0pt}
\setlength{\fboxrule}{2pt}

% Options: Sonny, Lenny, Glenn, Conny, Rejne, Bjarne, Bjornstrup
\usepackage[Conny]{fncychap}

% Redifining from 0.1 to 1 ... 2 ... and so on
\renewcommand{\thesection}{\arabic{section}}

% Renaming Bibliography name
\renewcommand\bibname{Bibliographie}

% Bibliography styling
\bibliographystyle{plain}

% Redifining from 1.1 to 1.a ... 1.b ... and so on
%\renewcommand{\thesubsection}{\alph{subsection}}

% Part rules
\titlecontents{part}[1em]
{\vskip 0.5ex}%
{}% numbered sections formatting
{\itshape\scshape}% unnumbered sections formatting
{}%

% Chapter rules
\titlecontents{chapter}[1em]
{\vskip 0.5ex}%
{}% numbered sections formatting
{\itshape\scshape}% unnumbered sections formatting
{\titlerule*[0.3pc]{.}\contentspage}%

% Section rules
\titlecontents{section}[1em]
{\vskip 0.2ex}%
{\hspace{0.15in}\contentslabel{0.13in}.\hspace{0.1in}}% numbered sections formatting
{\itshape}% unnumbered sections formatting
{\titlerule*[0.3pc]{.}\contentspage}%

% Subsection rules
\titlecontents{subsection}[3.5em]
{}%
{\contentslabel{0.10in}\hspace{0.25in}}% numbered sections formatting
{}% unnumbered sections formatting
{\titlerule*[0.3pc]{.}\contentspage}%

% Custom colors!
\definecolor{MYCOL}{HTML}{00318A}

% List of Figures settings:
\renewcommand{\listfigurename}{Liste Des Figures}
\renewcommand{\cftloftitlefont}{\vspace*{-0.6in}\hfill\color{Black}\fontsize{35}{43}\bfseries}
\renewcommand{\cftafterloftitle}{\hfill}
\setlength{\cftfigindent}{0pt} % left aligned Fig entries
\renewcommand{\cftfigpresnum}{Figure } % put before the number
\renewcommand{\cftfigaftersnum}{:} % put before the number
\addtolength{\cftfignumwidth}{2.5em} % extra space for \cftpresnum

% List of Tables settings:
\renewcommand{\listtablename}{Liste Des Tableaux}
\renewcommand{\cftlottitlefont}{\vspace*{-0.6in}\hfill\color{Black}\fontsize{35}{43}\bfseries}
\renewcommand{\cftafterlottitle}{\hfill}
\setlength{\cfttabindent}{0pt} % left aligned Tab entries
\renewcommand{\cfttabpresnum}{Tableau } % put before the number
\renewcommand{\cfttabaftersnum}{:} % put before the number
\addtolength{\cfttabnumwidth}{2.5em} % extra space for \cftpresnum

% Table of contents settings:
\renewcommand{\contentsname}{Table Des Matières}

\usepackage[framemethod=TikZ]{mdframed}

% For the box (theme)
\mdfdefinestyle{MyFrame}{%
    linecolor=black,
    outerlinewidth=2pt,
    roundcorner=0pt,
    innertopmargin=13pt,
    innerbottommargin=13pt,
    innerrightmargin=20pt,
    innerleftmargin=20pt,
    backgroundcolor=gray!10!white}


% My custom page style
\fancypagestyle{myfancy}{%   
   \fancyhf{}%
   \fancyfoot[C]{\tikz[baseline={(0,0)},anchor=center] \node[label={center:\thepage}]{\includegraphics[scale=.037]{../images/cloud.png}};}%
   \renewcommand{\headrulewidth}{0pt}%
   \renewcommand{\footrulewidth}{0pt}%
   \fancyhead{}
   \fancyfoot{}
   \fancyfoot[C]{\vspace{-0.35in}\textcolor{Gray}{\Large{Page}}\hspace{0.04in} \large{\thepage}\hspace{0.07in}\large{|}\hspace{0.07in}\large{\pageref{LastPage}}}
}%
\pagestyle{empty}

% Custom environment
\newenvironment{changemargin}[2]{%
\begin{list}{}{%
\setlength{\topsep}{0pt}%
\setlength{\leftmargin}{#1}%
\setlength{\rightmargin}{#2}%
\setlength{\listparindent}{\parindent}%
\setlength{\itemindent}{\parindent}%
\setlength{\parsep}{\parskip}%
}%
\item[]}{\end{list}}

% 'myfancy' page styling is also applied on chapters!
\usepackage{etoolbox}
\patchcmd{\chapter}{\thispagestyle{empty}}{\thispagestyle{myfancy}}{}{}

\titleformat{\chapter}
{\centering\fontsize{40}{50}\bfseries\color{Black}\scshape}
{}
{}
{\vspace{-1.7in}{\color{Black} \rule{\linewidth}{1.2mm} }}[\vspace{-0.35in}{\color{Black} \rule{\linewidth}{1.2mm} }\vspace{-1in}]

\titleformat{\section}
{\Large\bfseries\color{Black}}
{\thesection.\hspace{0.08in}}
{0em}
{}[\titlerule]

\titleformat{\subsection}
{\large\bfseries\color{Black}}
{\hspace{0.25in}\thesubsection\hspace{0.08in}}
{0em}
{}[]

\titleformat{\subsubsection}
{\bfseries\large}
{\hspace{0.40in}}
{0em}
{}[]

\newenvironment{myindentpar}[1]%
  {\begin{list}{}%
          {\setlength{\leftmargin}{#1}}%
          \item[]%
  }
  {\end{list}}

\usepackage{tikz}
\usetikzlibrary{calc}

\begin{document}

\begin{titlepage}

   \begin{tikzpicture}[remember picture, overlay]
      \draw[line width = 2pt] ($(current page.north west) + (0.25in,-0.25in)$) rectangle ($(current page.south east) + (-0.25in,0.25in)$);
   \end{tikzpicture}

   \begin{center}
       \vspace*{-0.4in}

       \begin{large}

           République Algérienne Démocratique et Populaire

           Ministère de l'Enseignement Supérieur et de la Recherche Scientifique

           Université M'Hamed Bougara de Boumerdès

       \end{large}

       \vspace{0.1in}

    \begin{figure}[h]
    \centering
        \includegraphics[width = 1.5in, height = 1in]{../images/UnivUMBB.jpg}
    \end{figure}
     \vspace*{-0.1in}
    \large{Faculté des Sciences}
    \\
    \vspace*{-0.1in}
    \large{Département d'informatique}
   \end{center}

    \hspace{0.2in}
    \large{\textbf{Domaine \,\; :}  Mathématiques Informatique}
    \kern 1in
    \large{\textbf{Année universitaire :}}

    \hspace{0.2in}
    \large{\textbf{Filière \quad\,\,\, :}  Informatique}
    \kern 2.65in
    2021 / 2022

    \hspace{0.2in}
    \large{\textbf{Spécialité \, :}  Technologie de l'information}

    \begin{center}
        \begin{large}

            \vspace{0.1in}

            \textit{\textbf{\uline{Projet fin d'étude}}}
        
        \end{large}

        \vspace{0.3in}

        \textit{\Huge{\textbf{\uline{Thème}}}}

        \vspace{0.2in}

        \begin{mdframed}[style=MyFrame]
            \begin{center}
            \color{Black}
              \Large{\textbf{La segmentation et Comptage des cellules biologiques}}

              \Large{\textbf{dans images de microscope par apprentissage profond}}
            \end{center}
        \end{mdframed}
    \end{center}

    \vspace{0.4in}

    \hspace{0.2in}
    \textit{\textbf{Présenté par :}}

    \hspace{0.2in}
    Gharbi Aghiles

    \hspace{0.2in}
    Neggazi Mohamed Lamine

\end{titlepage}

\newpage

\vspace*{0.2in}

\thispagestyle{empty}

\begin{center}
    {\color{Black} \rule{3in}{1.4mm} }\\
    \vspace{0.1in}
    \scshape{\fontsize{34}{46}{\bfseries{\color{Black}{Résumé}}}}
    \\
    \vspace{0.6in}
\end{center}
\cftaddtitleline{toc}{part}{\vspace{-0.12in}\color{Black}{Résumé}}{}
\begin{changemargin}{0.9cm}{0.9cm}
\hspace*{0.16in}
\end{changemargin}

Le comptage de cellules est une tâche fastidieuse qui bénéficierait grandement de l'automatisation. Le comptage précis des cellules fournit des informations quantitatives essentielles et joue un rôle clé dans la recherche biologique ainsi que dans les applications industrielles et biomédicales. Malheureusement, la méthode de comptage manuel couramment utilisée demande beaucoup de temps, mal standardisée et n'est pas reproductible. La tâche est rendue encore plus difficile par le chevauchement des cellules. la mauvaise qualité de l'imagerie ... , nous décrivons ici un réseau de neurones convolutionnel qui vas segmenter les cellules comme une première phase, et les compter dans la seconde phase.

\vspace{1in}

\begin{changemargin}{0.9cm}{0.9cm}
Mots clés : CNN, Comptage des cellules, segmentation d'image
\end{changemargin}

\newpage

\vspace*{0.2in}

\thispagestyle{empty}

\begin{center}
    {\color{Black} \rule{3in}{1.4mm} }\\
    \vspace{0.1in}
    \scshape{\fontsize{34}{46}{\bfseries{\color{Black}{Abstract}}}}
    \\
    \vspace{0.6in}
\end{center}
\cftaddtitleline{toc}{part}{\vspace{-0.12in}\color{Black}{Abstract}}{}
\begin{changemargin}{0.9cm}{0.9cm}
\hspace*{0.16in}

Cell counting is a tedious task that would benefit greatly from automation. Accurate cell counting provides essential quantitative information and plays a key role in biological research as well as industrial and biomedical applications. Unfortunately, the commonly used manual counting method is time consuming, poorly standardized and not reproducible. The task is made even more difficult by overlapping cells and poor imaging quality. In this paper we describe a convolutional neural network that will segment the cells as a first phase, and count them in the second phase.

\end{changemargin}

\vspace{1in}

\begin{changemargin}{0.9cm}{0.9cm}
Keywords: CNN, cell segmentation, cell counting, convolutional neural networks
\end{changemargin}

\newpage

\cftaddtitleline{toc}{part}{\vspace{-0.12in}\color{Black}{Table des matières}}{}
\tableofcontents

\newpage

\cftaddtitleline{toc}{part}{\vspace{-0.12in}\color{Black}{Liste des figures}}{}
\listoffigures

\newpage

\cftaddtitleline{toc}{part}{\vspace{-0.12in}\color{Black}{Liste des tableaux}}{}
\listoftables

\addtocontents{toc}{\protect\renewcommand{\protect\cftsecleader}{\protect\cftdotfill{\protect\cftsecdotsep}}}

\addtocontents{lof}{\protect\thispagestyle{empty}}
\addtocontents{lot}{\protect\thispagestyle{empty}}
\addtocontents{toc}{\protect\thispagestyle{empty}}

\newpage

\pagestyle{myfancy}

\vspace*{-0.2in}

\setcounter{page}{1}

\begin{center}
    {\color{Black} \rule{6.2in}{1.4mm} }\\
    \vspace{0.1in}
    \scshape{\fontsize{34}{46}{\bfseries{\color{Black}{Introduction générale}}}}
    \\
    \vspace{0.5in}
\end{center}
\addcontentsline{toc}{chapter}{\vspace{-0.12in}\color{Black}{Introduction générale}}
\hspace*{0.16in}

\newpage

\vspace*{\fill}
\begin{center}
    {\color{Black} \rule{\linewidth}{1.2mm} }\\
\vspace{0.25in}
 {\centering\fontsize{30}{40}{\bfseries{\color{Black}{\scshape{Chapitre I : Etat de l'art}}}}}
\vspace{0.35in}\\
    {\color{Black} \rule{\linewidth}{1.2mm} }
\end{center}
\vspace*{\fill}
\addcontentsline{toc}{chapter}{\color{Black}{Chapitre I : Etat de l'art}}
\setcounter{section}{0}

\newpage

\section{Introduction}
\vspace{0.2in}
\hspace*{0.16in}

\section{Related Work}
\vspace{0.2in}
\hspace*{0.16in}
Blood cell segmentation is the extraction of different blood cells from microscopic images. Blood cell counting is the process of counting the detected blood cells after segmentation. Many researchers have implemented methods for segmentation and counting of blood cells.\\

V.KIMBAHUNE and al \textsuperscript{\cite{V.KIMBAHUNE}} have developed a method for segmenting and counting red blood cells (RBC) and white blood cells (WBC).
segmentation is done using Pulse-Coupled Neural Network (PCNN) and square tracing algorithm for countour tracing after de-noising it with PCNN + median filter, the counting is performed by scaning the image and using edge detection method. this method gave good results compared to state of art methods.

Bhavnani and al \textsuperscript{\cite{bhavnani2016segmentation}} have developed a method for segmenting and counting RBC (red blood cells), WBC (White blood cells) and platelets which is also called complete blood count (CBC), by using Otsu’s thresholding and morphological operations as a segmentation method, and for counting the are preforming a comparison between two methods: the watershed algorithm and Circular Hough Transform. the CHT method is the best in terms of accuracy but it has some weaknesses with overlaping cells and morphological abnormalities.

Carlos X. Hern{\'{a}}ndez et al \textsuperscript{\cite{DBLP:journals/corr/abs-1802-10548}} have implemented a convolutional neural network (CNN) using a feature pyramid network combined with a VGG style neural network for segmenting and counting of cells in a given microscopy image which achieved a relatively good accuracy of 95\% but with some failure cases such as: High cell overlap, Irregular cell shapes and bad focal planes.

K. Sudha and P. Geetha \textsuperscript{\cite{SUDHA2020639}} have developed a two stage framework which will segment the leukocytes (a type of WBC) with an edge strength-based Grabcut method as a first stage, in the second stage will count the cells using the novel gradient circular hough transform (GCHT) method. the new GCHT method can segment touched cells and even overlapped cells.

Overton, Toyah and Tucker, Allan \textsuperscript{\cite{10.1007/978-3-030-44584-3_31}} have developed method about segmentation and counting IDP(Internally Displaced people) and erythrocytes (red blood cells) using DO-U-Net (Dual Output U-Net) which outputs a segmentation mask and an edge mask then they substract them to get rid of the overlapping and the touching problem, the model trains on extremely small datasets (10 images) and gives a high segmentation accuracy, for the IDP they had 98.69\% for fixed resolution images and 94.66\% for scale-invariant satellite images.

\section{Conclusion}
\vspace{0.1in}
\hspace*{0.16in}


\newpage

\vspace*{\fill}
\begin{center}
    {\color{Black} \rule{\linewidth}{1.2mm} }\\
\vspace{0.25in}
    {\centering\fontsize{30}{40}{\bfseries{\color{Black}{\scshape{Chapitre II : Analyse et Conception}}}}}
\vspace{0.35in}\\
    {\color{Black} \rule{\linewidth}{1.2mm} }
\end{center}
\vspace*{\fill}
\addcontentsline{toc}{chapter}{\color{Black}{Chapitre II : Analyse et Conception}}
\setcounter{section}{0}

\newpage

\section{Introduction}
\hspace*{0.16in}

\section{Conclusion}
\hspace*{0.16in}

\newpage

\vspace*{\fill}
\begin{center}
    {\color{Black} \rule{\linewidth}{1.2mm} }\\
\vspace{0.25in}
    {\centering\fontsize{30}{40}{\bfseries{\color{Black}{\scshape{Chapitre III : Implémentation}}}}}
\vspace{0.35in}\\
    {\color{Black} \rule{\linewidth}{1.2mm} }
\end{center}
\vspace*{\fill}
\addcontentsline{toc}{chapter}{\color{Black}{Chapitre III : Implémentation}}
\setcounter{section}{0}

\newpage

\section{Introduction}
\vspace{0.2in}
\hspace*{0.16in}

\section{Conclusion}
\vspace{0.1in}
\hspace*{0.16in}

\newpage

\vspace*{-0.2in}

\begin{center}
    {\color{Black} \rule{5.5in}{1.4mm} }\\
    \vspace{0.1in}
    \scshape{\fontsize{34}{46}{\bfseries{\color{Black}{Conclusion générale}}}}
    \\
    \vspace{0.5in}
\end{center}
\addcontentsline{toc}{chapter}{\vspace{-0.12in}\color{Black}{Conclusion générale}}
\hspace*{0.16in}

\newpage

\bibliography{main}

\end{document}
